% 本文件由 LaTeX 排版 Agent 根据有效 translation.json 自动生成。
% author=Councils; work_id=3803; title=新凯撒勒雅会议（公元315年）

\chapter{\WorkTitle{新凯撒勒雅会议（公元315年）}}

% structure: p0001 source=h1 target=omitted reason=page-title-already-rendered-by-parent

会议聚集于新凯撒勒雅，即本都的一座城市，在顺序上仅次于安西拉会议，而日期早于其他会议，甚至早于在尼西亚举行的第一届大公会议。在此会议中，圣父们聚集在一起，其中有圣殉道者巴西略（阿马塞亚主教），采纳了如下规条以建立教会秩序——\par

圣洁可敬的教父们在新凯撒勒雅聚会的规条，这些规条确实日期晚于安西拉会议所制定的，但早于尼西亚会议：然而，因其特殊尊严，尼西亚会议被置于它们之前。\par

% structure: p0004 source=h2 target=section reason=relative-heading-rank; base=h2

\section{规条}

% structure: p0005 source=h3 target=subsection reason=relative-heading-rank; base=h2

\subsection{第一条}

若一司铎结婚，应将其从神职中撤除；但若他犯私通或奸淫，则应完全被逐出（即，共融）并使其补赎。\par

% structure: p0007 source=h3 target=subsection reason=relative-heading-rank; base=h2

\subsection{第二条}

若一妇人嫁给了两兄弟，应被逐出（即，共融）直至其死亡。然而，临终时，作为仁慈之举，她可被接纳补赎，前提是她声明若康复将解除婚姻。但若此婚姻中的妇人或男子死亡，幸存者的补赎将非常困难。\par

% structure: p0009 source=h3 target=subsection reason=relative-heading-rank; base=h2

\subsection{第三条}

关于那些多次结婚者，规定的补赎时间是众所周知的；但他们的生活方式和信德可缩短时间。\par

% structure: p0011 source=h3 target=subsection reason=relative-heading-rank; base=h2

\subsection{第四条}

若有人贪恋一女子并计划与她同寝，但计划未实现，则显然他赖恩宠得救。\par

% structure: p0013 source=h3 target=subsection reason=relative-heading-rank; base=h2

\subsection{第五条}

若一望教者进入教会并已列于望教者等级中，而犯罪，则他若为跪拜者，应成为听众且不再犯罪。但若他在作听众时再次犯罪，则应被逐出。\par

\EditorialNote{若望·佐纳拉斯：\Quote{有两种望教者。有些是刚进来的，这些人尚不完全，在圣经和福音诵读后立即出去。但另有一些人已准备一段时间并获得一定成全；他们在福音后等待为望教者的祈祷，当听到\InnerQuote{望教者，向主低头}时，他们跪下。这些人，作为更成全的，已尝过天主的善言，若跌倒，则被移除原位，被置于听众中；但若有人在作听众时犯罪，则完全被逐出教会。}}\par

% structure: p0016 source=h3 target=subsection reason=relative-heading-rank; base=h2

\subsection{第六条}

关于怀孕妇人，规定她可随时愿意领洗；因为在此事上，妇人并不传递给婴儿任何东西，因为公开宣认信仰显然是每个人的个人[特权]。\par

% structure: p0018 source=h3 target=subsection reason=relative-heading-rank; base=h2

\subsection{第七条}

\Emph{司铎}不应出席缔结第二次婚姻者的婚礼；因为既然再婚者应受补赎，那么出席婚宴并认可此婚姻的司铎，又该是怎样的人呢？\par

% structure: p0020 source=h3 target=subsection reason=relative-heading-rank; base=h2

\subsection{第八条}

若一平信徒的妻子犯奸淫并被明确证实，则该[丈夫]不能进入圣职；若她在其晋秩后犯奸淫，他必须休妻；但若他留住她，则他不能在委托给他的圣职中有份。\par

% structure: p0022 source=h3 target=subsection reason=relative-heading-rank; base=h2

\subsection{第九条}

\Emph{司铎}在犯肉身之罪后被提升，且他告解说在晋秩前已犯罪，则不得举行祭献，但因他在其他方面的热忱，仍可保留其他职务；因为多数人确认晋秩可涂去其他种类的罪。但若他不告解且不能被公开定罪，决定取决于他自己。\par

% structure: p0024 source=h3 target=subsection reason=relative-heading-rank; base=h2

\subsection{第十条}

同样，若一执事陷入同一罪行，让他担任服务者的等级。\par

% structure: p0026 source=h3 target=subsection reason=relative-heading-rank; base=h2

\subsection{第十一条}

司铎不应在三十岁前晋秩，即使他在各方面都是配得的人，也应让他等待。因为我们的主耶稣基督在三十岁时受洗并开始教导。\par

% structure: p0028 source=h3 target=subsection reason=relative-heading-rank; base=h2

\subsection{第十二条}

若有人在患病时领洗，因为其（宣认）信仰并非自愿而是出于必需（即，因怕死），他不可被提升至司铎职位，除非因其后来的（显示出的）热忱和信德，并因缺乏人手。\par

\EditorialNote{希腊文用于\Quote{受洗}的词是\Quote{被光照}（\OriginalTerm{被光照}{\GreekText{φωτισθῇ}}），这是古人中非常常见的表达。}\par

\EditorialNote{佐纳拉斯解释此禁令的理由是众所周知的事实：在那个时代，人们推迟领洗，以便更长时间地不受洗礼被认为施加的约束。}\par

% structure: p0032 source=h3 target=subsection reason=relative-heading-rank; base=h2

\subsection{第十三条}

当城市主教或司铎在场时，乡村司铎不得在城市的教会中举行祭献；也不得在祈祷时分发饼或杯。然而，若他们缺席，而他（即，一乡村司铎）被单独召请祈祷，他可分发。\par

\EditorialNote{许多学者，包括狄奥尼修斯·埃克西古和依西多尔，将最后一款读作复数。在许多手稿中，此条规与下一条合并，总数记为14。}\par

% structure: p0035 source=h3 target=subsection reason=relative-heading-rank; base=h2

\subsection{第十四条}

然而，\OriginalTerm{乡村主教}{chorepiscopi}确实仿效七十门徒的模式；作为同仆，因其对穷人的奉献，享有举行祭献的荣誉。\par

\EditorialNote{见上一条规下的注释，因为在许多手稿中，两条规在古代摘要中是合并的。}\par

% structure: p0038 source=h3 target=subsection reason=relative-heading-rank; base=h2

\subsection{第十五条}

根据规条，执事应为七位，即使城市很大。对此，你将从\BibleBook{宗徒}{大事录}中确信。\par

这一规条在罗马被遵守，直到十一世纪，七位枢机执事的数目才改为十四位。格拉先将它收入\WorkTitle{法令集}（第一部分，第XCIII区别，第xii章）是很好的证据，表明他认为这是罗马纪律的一部分。欧瑟比给出了教宗科尔内略的一封信，写于第三世纪中期，其中说当时在罗马有四十四位司铎、七位执事和七位副执事；较低等级的人数非常多。托马西努斯说：\Quote{无疑，罗马教会在此事上意图效法宗徒，他们只祝圣了七位执事。但其他教会并不那样严格地保持那个数目。}\par

在加采东大公会议的法案中，记载埃德萨教会有十五位司铎和三十八位执事。我们知道，犹斯定为君士坦丁堡教会任命了一百位执事。凡·埃斯彭很好地指出，虽然此规条提及先前关于该主题的法律，但会议本身以及希腊注释家巴尔萨蒙或佐纳拉斯都没有给出丝毫提示该规条是什么。\par

新凯撒勒雅的教父们将一城中执事数目限于七位是基于圣经的权威，但\OriginalTerm{五-六会议}{Quinisext}会议的第十六条规明确说，他们这样做表明他们指的是赈济职员，而非神圣奥迹职员，并且圣斯德望和其他人完全不是后一种意义上的执事。读者可参考此规条，在那里为了捍卫君士坦丁堡的实践，我们所考虑的这条规条的含义被完全歪曲了。\par

\SourceRecord{Translated by Henry Percival.；Nicene and Post-Nicene Fathers, Second Series；Vol. 14.；Edited by Philip Schaff and Henry Wace.；Buffalo, NY: Christian Literature Publishing Co.,；1900.；<http://www.newadvent.org/fathers/3803.htm>.}
